%%=============================================================================
%% Projectgegevens - VUL DIT ALS EERSTE IN
%%=============================================================================
%%
%% Dit is de enige plek waar je je naam, je titel, je bedrijf en je begeleiders
%% invult. Het voorstel, het rapport en de poster halen hun gegevens hier op,
%% zodat ze nooit uit elkaar kunnen lopen.
%%
%% Enkel de tekst tussen de accolades aanpassen. Laat de commandonamen staan.

%% ---------- Jij ------------------------------------------------------------
\newcommand{\studentnaam}{Xiaohong Zhou}
\newcommand{\studentemail}{xiaohong.zhou@student.hogent.be}

%% Wordt gebruikt als bestandsnaam van je PDF's, bijvoorbeeld
%% Familienaam-Voornaam-rapport.pdf. Geen spaties of leestekens gebruiken.
\newcommand{\bestandsnaam}{Zhou-Xiaohong}

%% ---------- Je project -----------------------------------------------------
\newcommand{\projecttitel}{Ontwikkelen webpagina voor automatisering branddetectieverslag}

%% Een ondertitel is optioneel. Laat leeg als je er geen hebt: {}
\newcommand{\projectondertitel}{}

%% Een of twee zinnen. Wordt gebruikt op de poster en in het voorstel.
\newcommand{\projectpitch}{%
    Dit project maakt gebruik van PHP/Laravel, MySQL en AWS en zorgt voor een koppeling tussen databank en formulier zodat het niet steeds van nul moeten worden opgesteld.%
}

%% ---------- Waar je het doet ------------------------------------------------
\newcommand{\bedrijfsnaam}{Belgateq}

%% Je mentor op de werkvloer (het bedrijf)
\newcommand{\mentornaam}{Burak Yorulmaz}
\newcommand{\mentoremail}{it@belgateq.be}

%% Je begeleider van de opleiding (HOGENT)
\newcommand{\begeleidernaam}{Isabelle Demeulenaere}
\newcommand{\begeleideremail}{isabelle.demeulenaere@hogent.be}

%% ---------- Je repository ---------------------------------------------------
%% De link naar je eigen git-repository. Die hoort zichtbaar in je rapport en
%% op je poster: hij is een beoordeeld onderdeel van je oplevering.
\newcommand{\repolink}{https://github.com/xiaohongzhou8/mijn-graduaatsproject}

%% ---------- Vast voor dit opleidingsonderdeel -------------------------------
%% Deze twee hoef je normaal niet aan te passen.
\newcommand{\opleidingnaam}{Graduaat in het Programmeren}
\newcommand{\academiejaarwaarde}{2026-2027}

%% Zoekwoorden voor het voorstel en de poster (3 tot 5)
\newcommand{\projectkeywords}{trefwoord, trefwoord, trefwoord}
